% !TEX program = latexmk (xelatex)
\documentclass[12pt, a4paper]{article}

\usepackage[polish]{babel}
\usepackage{fontspec}
\usepackage{lato}
\renewcommand{\familydefault}{\sfdefault}
\setlength{\parskip}{0.3em}
\usepackage{fancyhdr}

\begin{document}
\pagestyle{myheadings}
\markright{PZSP2 konsultacje z oprogramowania 1 \hfill \today \hfill}

\section{Cel projektu}
Chcemy stworzyć system rezerwacji dla szkoły tańca. Będzie to strona internetowa
z interfejsem dla: uczniów, instruktorów i dla właścicielki.

\section{Technologie}
Pracując nad projektem, używamy konteneryzacji. Pozwala to ujednolicić środowisko
deweloperskie oraz zapobiec niespodziankom podczas wdrożenia.

Frontend piszemy w Typescrip'cie, używając frameworku React i narzędzia
do budowania Vite. Tak wybraliśmy, bo znamy React'a z poprzednich projektów.

Backend to Django w Pythonie - jeden z prostszych sposobów na napisanie
backendu webowego., prostszy od np. Springa. Plus, Ihor ma z nim doświadczenie.

Baza danych to PostgreSQL. Otrzymaliśmy ciekawą sugestię użycia SQLite
-- wtedy oszczędzilibyśmy na jednym kontenerze, co uprościłoby architekturę,
lecz zrobiło ją mniej elastyczną. Gdy baza jest oddzielnie, można w przyszłości
dodać jakiś mikroserwis. Nie planujemy tego robić, ale architekturę chcemy mieć
„pełną” -- pozostali eksperci przychylniej na to patrzą.

Potok CI/CD składa się z budowania frontendu, lintowania i testowania frontendu oraz
testowania backendu. Są to czynności podstawowe, które każdy programista wykonuje
na bieżąco podczas pracy z kodem, więc z odpalania potoku „na gałęzi”, korzyść
dla programisty jest nikła. Dlatego, potok można włączyć ręcznie przy otwartym
merge reques'cie. Pomyślne wykonanie potoku jest warunkiem koniecznym
do przyjęcia MR'a.

\section{Higiena pracy}
Na gitlabie tworzymy isia. Programista tworzy draft MR na to isiu i implementuje.
Jak uzna, że gotowe, to odznacza „draft” i MR czeka na recenzję \textit{jednej} osoby.
Jednym z warunków wdrożenia jest, że gałąź musi być przedłużeniem gałęzi
głównej (rebase).

\section{Testy akceptacyjne}
Chcąc zasymulować scenariusz przemysłowego wytwarzania oprogramowania,
opracowaliśmy schemat testów akceptacyjnych.

\subsection{Rejestracja na zajęcia}
Uczeń, po założeniu konta, będzie mógł przejrzeć dostępne grupy zajęciowe
i~zapisać się na wybraną. Gdy przebiegnie to pomyślnie, zostanie o tym poinformowany,
będzie figurował na liście zapisanych oraz odpowiednio zwiększy się kwota jego karnetu.

\subsection{Odrabianie nieobecności}
Jeśli na przewidzianych dla danego ucznia zajęciach zostanie odnotowana
jego nieobecność,
na stronie pojawi się informacja o możliwości jej odrobienia.
Na dowolnych innych zajęciach, instruktor będzie miał możliwość dopisania tej osoby
do standardowej listy obecności, co „odrobi” nieobecność.

\subsection{Płatności}
Jeśli uczeń nie dokonał jeszcze płatności za obecny miesiąc lub któryś poprzedni,
na jego koncie powinna widnieć o tym informacja. Właścicielka powinna widzieć
taką samą informację oraz mieć możliwość ręcznego oznaczenia tej płatności
jako uregulowanej.

\raggedleft
\vspace{1em}
\begin{tabular}{l}
    Zespół 2 \\\hline
    \rule{0pt}{1.5em}Gavrylenko Ihor\\
    Kukiełka Wojciech\\
    Rybachok Andrii\\
    Yedzeika Sofiya
\end{tabular}

\end{document}
